\documentclass[10pt]{article}
\usepackage[margin=20mm]{geometry}
\usepackage{amsmath,amssymb,mathtools,fontspec,unicode-math,xurl,hyperref}
\IfFontExistsTF{DejaVu Serif}{\setmainfont{DejaVu Serif}}{\setmainfont{Latin Modern Roman}}
\setmathfont{latinmodern-math.otf}
\setlength{\parindent}{0pt}\setlength{\parskip}{6pt}\sloppy
\title{Initial Technical Review of M25-0411}
\author{AI-assisted technical review}
\date{}
\begin{document}
\maketitle
Abstract: This review examines the classification of equivariant harmonic map heat flows and bounded radial critical heat flows in the specified preprint of Kim and Merle. No substantive issue was established in the inspected scope. The review checks the theorem hypotheses, the dependence on soliton resolution, the modified-profile and dissipation arguments, and the transition from modulation estimates to the stated rates. Independent calculations verify the signs and normalizations of the interaction constants and the rate recurrence. This is an initial technical report; scientific adjudication remains pending.\par
\section*{Paper identity and reviewed scope}
The paper is Kihyun Kim and Frank Merle, On classification of global dynamics for energy-critical equivariant harmonic map heat flows and radial nonlinear heat equation, associated with Communications on Pure and Applied Mathematics 78 (2025), 1783–1842, DOI https://doi.org/10.1002/cpa.22253. The reviewed object is the 46-page PDF of arXiv:2404.04247v3, submitted 1 April 2025 [KM]. All page locators below refer to that PDF, whose printed page numbers agree with the PDF page numbers. The journal version was not compared line by line; this report does not silently transfer its conclusion to an uninspected version.\par
The complete supplied PDF, including Appendix A and its bibliography, was read. All 46 source pages were also visually inspected through the native image route. The main focus was Theorems 1.1, 1.2 and 1.4; Lemmas 2.6–2.7; Proposition 3.1 and its construction; Proposition 4.1; and Proposition 5.2 and Section 5.2. External theorems were not all independently reproved. In particular, the hypotheses of Aryan’s soliton-resolution theorem were checked in the explicitly identified later version [Aryan], pp. 1–3; the full proof of that external paper was not reviewed.\par
The frozen manifest and a bounded check of the current arXiv entry, author publication page and publisher metadata did not identify a separate correction to this reviewed version. This is a statement about the search performed, not a proof that no correction exists. No relevant separate executable supplement was located in the acquired manifest. No author code was executed, and absence of code is not treated as a mathematical defect.\par
\section*{Definitions and operative assumptions}
For harmonic map heat flow (HMHF), the unknown is the real radial angle \(v(t,r)\) with equivariance index \(D\in\mathbb Z_{\geq3}\): \[\partial_t v=\partial_{rr}v+r^{-1}\partial_rv
             -\frac{D^2}{2r^2}\sin(2v).\] Theorem 1.1, p. 4, concerns every finite-energy datum in the sector \(\mathcal E_{\ell,m}\), with endpoint limits \(\ell\pi\) and \(m\pi\), where \(\ell,m\in\mathbb Z\). It asserts global existence, \(J=|m-\ell|\) bubbles, a common sign \(\operatorname{sgn}(m-\ell)\) when \(J>0\), and the stated universal scales. The case \(J=0\) is the empty-sum, decaying case; the unused limiting scale then has no uniqueness content.\par
For the nonlinear heat equation (NLH), the statement is restricted to radial solutions of \[\partial_tu=\Delta u+|u|^{4/(N-2)}u,\qquad N\geq7,
 \qquad \sup_{0\leq t<T}\|u(t)\|_{\dot H^1}<\infty.\] Theorem 1.2, p. 5, expressly assumes the statement formulated as Open Question 2.2, p. 13. It therefore does not assert the corresponding classification for arbitrary unbounded, nonradial, or lower-dimensional solutions. Remark 1.3, p. 5, and the paragraph following (2.4), p. 13, explicitly acknowledge the subsequent soliton-resolution work.\par
The common notation is \[D=(N-2)/2>2,\qquad p=1+2/D,\qquad
 W(r)=\left(1+\frac{r^2}{N(N-2)}\right)^{-D}\] for NLH. For HMHF, \(u=r^{-D}v\), \(N=2D+2\), \(Q(r)=2\arctan(r^D)\) and \(W=r^{-D}Q\). The respective scaling generators are \(\Lambda W=(r\partial_r+D)W\) and \(\Lambda Q=r\partial_rQ\). The transformed HMHF states belong to affine sectors \(\mathcal H_{0,m}\). The paper requires the remainder after subtracting bubbles and the body map to lie in \(\dot H^1\); it does not require the individual transformed HMHF bubble \(r^{-D}Q\) to lie in that space.\par
Theorem 1.4, p. 10, fixes a finite integer \(J\geq1\), fixed signs, continuous positive scales satisfying \(\lambda_J\ll\cdots\ll\lambda_1\), and convergence of the remainder in \(\dot H^1\). A finite endpoint requires \(\lambda_1(t)\to0\) and an appropriate body map \(z^*\) in the stated affine sector; the global case has \(z^*=0\). These conditions, including the ordered scale ratios, are used essentially. The proof subsequently chooses differentiable modulation parameters asymptotic to the original continuous scales; it does not assume that the initially supplied scales were differentiable.\par
\section*{Proof dependencies and checks}
\subsection*{The external resolution hypothesis}
The reduction on pp. 9–13 uses Proposition 2.1 for HMHF and the statement of Open Question 2.2 for NLH. In the later reference actually inspected, Aryan’s Theorem 1.1 [Aryan], p. 2, assumes radial finite-energy data and boundedness of the energy-space norm near the maximal endpoint. It gives continuous, separated bubble scales, an energy-small remainder, and a body term only at a finite endpoint.\par
This matches the required form after reversing the scale order and renaming the dimension. Indeed, Hardy’s inequality for \(N\geq3\) gives \[\int_0^\infty |u|^2r^{N-3}\,dr
 \leq \frac{4}{(N-2)^2}
       \int_0^\infty |\partial_ru|^2r^{N-1}\,dr.\] Thus the norm used there, which also includes the Hardy term, is equivalent to the radial \(\dot H^1\) norm here. The finite-endpoint separation from \(\sqrt{T-t}\) implies \(\lambda_1\to0\). The displayed conditional statement in the reviewed paper is consequently not itself an unsupported unconditional assertion. Checking the later theorem’s hypotheses is not a retrospective certification of every argument in its 2024 predecessor.\par
\subsection*{Modified profiles and dissipation}
The essential chain is \[\text{Lemma 2.7}\ \longrightarrow\ \text{Proposition 3.1}\
 \longrightarrow\ \text{Proposition 4.1}\
 \longrightarrow\ \text{Proposition 5.2 and Section 5.2}.\] Lemma 2.7, pp. 15–16, is a norm coercivity statement for the linearized operator after removing one kernel direction at each scale. It is not a claim that the NLH quadratic form has no negative directions. The compactness argument uses fixed finite \(J\), separated scales, orthogonality to a compactly supported test function with nonzero kernel pairing, and Hardy–Rellich bounds in the stated dimension range.\par
The modified profiles of Proposition 3.1, pp. 16–17, are constructed in pp. 18–29 and Appendix A, pp. 42–44. Their purpose is to make the leading interaction nearly orthogonal to the linearized remainder. The order of the small parameters in (3.26), the local support of the correction, and the Hölder exponent \(p-1>0\) in the derivative iteration are material. In Lemma 3.6, p. 24, the derivative differences are summable because one can choose \(2^{-(p-1)}<a<1\). This remains possible for every fixed \(D>2\), without asserting a uniform contraction rate as \(D\to\infty\).\par
Equation reference: eq:diss. The weighted remainder estimates (3.10)–(3.12) have different roles: the small \(L^2\) remainder yields the dissipation estimate, while the weighted \(L^1\) estimates control the modulation pairings. Proposition 4.1, pp. 29–35, proves \[
 \int_{t_0}^{T}\left(\mathcal D+\|g\|_{\dot H^2}^2\right)\,dt<\infty,
 \qquad
 \mathcal D=\sum_{j=2}^J\frac{\mu_j^{2D}}{\lambda_j^2},
 \quad \mu_j=\frac{\lambda_j}{\lambda_{j-1}}.\] For finite \(T\), the proof first localizes near the origin to make the body map small, then uses Lemma 4.5 for the exterior. For infinite \(T\), global smallness of \(g\) supplies the simpler argument. Treating a finite-endpoint body map as globally small would omit this distinction.\par
\subsection*{Why the dimension threshold is operative}
For NLH, \(\Lambda W(r)\) has a nonzero leading term proportional to \(r^{-2D}\) at infinity. Hence the radial integrand for \(\|r\Lambda W\|_{L^2(\mathbb R^{2D+2})}^2\) behaves like \(r^{3-2D}\). It is integrable precisely when \(D>2\); at \(D=2\) the tail diverges logarithmically. The origin is harmless. This verifies the dimension restriction invoked on p. 10. The identity used on pp. 39–40, \(\Lambda_{-1}\Lambda W=(2-D)\Lambda W+R\), with faster-decaying \(R\), also needs \(2-D\ne0\). The lower-dimensional regimes discussed on p. 6 therefore do not furnish counterexamples to the stated theorems.\par
\subsection*{From modulation estimates to rates}
Proposition 5.2, pp. 37–40, contains a derivative of an \(o(1)\) correction as well as errors. It cannot be replaced by an exact unperturbed ODE without justification. Section 5.2, pp. 40–42, removes those errors by using \(0<s_2<\cdots<s_J<2\), sign-definite correction terms, and the integrable quantity in [eq:diss] (\#eq:diss). It obtains \(\lambda_k^\pm=e^{o(1)}\lambda_k\) with two-sided differential inequalities.\par
For \(k=1\), \(\mu_1=0\). The upper and lower logarithmic quantities are monotone in opposite directions and differ by \(o(1)\). Each bounds the other near the endpoint, so their common limit is finite, rather than an extended infinite value. This yields \(0<L_\infty<\infty\) and excludes finite \(T\). For \(k\geq2\), integrating the reciprocal power \((\lambda_k^-)^{-(D-2)}\) forces \(\kappa\,\iota_k\iota_{k-1}<0\), since otherwise its positive left side would eventually be bounded above by a negative divergent term. This is the decisive sign argument, not an assumption about the signs.\par
\section*{Independent calculations}
\subsection*{Interaction constants}
The definitions (1.11)–(1.12), p. 9, give opposite signs in the two equations. These signs control the final alternating NLH pattern and the common HMHF sign.\par
Claim 1.  For the HMHF normalization of the reviewed paper, \[\kappa_{\rm HMHF}=-\frac{4D^2}{\pi}\sin\frac{\pi}{D}<0
 \qquad(D\geq3).\]\par
Proof. Using \(\Lambda Q=2Dr^D/(1+r^{2D})\) and the substitution \(s=r^{2D}\), \[\int_0^\infty(\Lambda Q)^3r^{D-1}\,dr
 =4D^2\int_0^\infty\frac{s}{(1+s)^3}\,ds=2D^2.\] The same substitution and the beta integral give \[\int_0^\infty(\Lambda Q)^2r\,dr
 =2D\,B(1+1/D,1-1/D)=\frac{2\pi}{\sin(\pi/D)}.\] Substitution into (1.11) proves the formula. ◻\par
For NLH put \(c=N(N-2)=4D(D+1)\) and \(s=r^2/c\). Then \(\Lambda W=D(1-s)(1+s)^{-D-1}\). Direct beta integration gives \[\int W^pr^{N-1}\,dr=\frac{c^{D+1}}{2(D+1)},\qquad
 \int(\Lambda W)^2r^{N-1}\,dr
 =c^{D+1}\frac{D(D+2)}{D-1}B(D+1,D+1).\] Consequently, \[\kappa_{\rm NLH}=
 \frac{D-1}{2(D+1)(D+2)B(D+1,D+1)}>0.\] The angular surface factor cancels in each ratio.\par
Independent standard-library quadrature, performed without author code, corroborated these identities: for HMHF at \(D=3,4,5\), the absolute differences were below \(4\times10^{-15}\); for NLH at \(D=5/2,3,4,5\), they were below \(10^{-11}\). The analytic calculation above establishes the formulas; the finite quadrature is not an interval-certified universal proof.\par
\subsection*{Rate and scaling recurrence}
After the sign is fixed, the leading scale equation has the form \[\frac{\lambda_j'}{\lambda_j}
   =-|\kappa|\frac{\lambda_j^{D-2}}{\lambda_{j-1}^{D}}.\] Substitution of \(\lambda_j=L^{1+2\alpha_j}\beta_jt^{-\alpha_j}\) requires \[(D-2)\alpha_j=D\alpha_{j-1}+1,\qquad
 \beta_j=\left(\frac{\alpha_j}{|\kappa|}\right)^{1/(D-2)}
          \beta_{j-1}^{D/(D-2)}.\] With \(\alpha_1=0\) and \(\beta_1=1\), this yields exactly (1.9)–(1.10). The power of \(L\) is consistent because \((D-2)(1+2\alpha_j)-D(1+2\alpha_{j-1})=0\). Exact rational arithmetic independently checked these identities through \(j=5\) for \(D=5/2,3,4,5\). For example, the first four exponents are \((0,2,12,62)\) for \(N=7\) and \((0,1,4,13)\) for \(D=3\). This check verifies the rate algebra conditional on the rigorously controlled modulation estimates; it is not a replacement for those estimates.\par
\section*{Disposition and limitations}
No substantive Problem is asserted. Several apparent discrepancies were tested against the actual PDF and their downstream use. In particular, (3.25), p. 20, contains \(\mu_{J+1}\) inside the parenthesis, not \(\mu_{J+1}^D\); the image resolves a misleading extraction. The reversed dissipative sign in introductory displays on pp. 2–3 is a typographical defect: the proof uses the correct energy decrease explicitly on pp. 10 and 31. It does not establish an invalid core derivation or a false final theorem and is not elevated to an in-scope mathematical finding.\par
The complete supplied proof was read, but the review does not claim a machine-checked formal proof, a new PDE simulation, or a full independent reproof of every cited theorem. The bound is the exact reviewed version, the dependency and assumption checks described above, and the retained independent calculations. No correction to the main result is proposed.\par
Dimension | Recorded state\par
Initial technical result | No Issue Found within inspected scope\par
Scientific adjudication | Pending\par
Independent acceptance | Not performed by this report\par
Human reviews | None recorded\par
Author adoption of a correction | Not applicable; no correction proposed\par
9 K. Kim and F. Merle, On classification of global dynamics for energy-critical equivariant harmonic map heat flows and radial nonlinear heat equation, arXiv:2404.04247v3, 1 April 2025. https://arxiv.org/abs/2404.04247v3. Journal identity: https://doi.org/10.1002/cpa.22253. S. Aryan, Soliton resolution for the energy-critical nonlinear heat equation in the radial case, arXiv:2405.06005v3, 19 February 2026. Consulted pp. 1–3, especially Theorem 1.1 on p. 2. https://arxiv.org/abs/2405.06005v3. K. Kim, publication list, consulted 29 October 2026. https://sites.google.com/view/kihyunkim.\par
\section*{Supporting checks}
The manifest SHA256 is `a99340df081a6a3d04d3be07dda54f76fd17dfcf7ed3b4625f7ffc90a619ece1`. The separately acquired Aryan PDF has SHA256 `cb14d0e3260cadf6c281b3e98f95f9aefac2c62471c8bcc802c3ba4a9dd6e50b`.\par
The frozen template has SHA256 `e450ad2d76f5d0dcf5585a948e78a31ed44c994ad6f3421ba54478afe63d1e37`.\par
The scientific analysis and actual image observations were performed by the primary executor in this conversation. A native helper performed only clerical source binding and text extraction. No human mathematical review is claimed.\par
\end{document}
